% !TeX spellcheck = en_US
\documentclass[french]{yLectureNote}

\title{Algèbre linéaire 3*}
\subtitle{Physique}
\author{Paulhenry Saux}
\date{\today}
\yLanguage{Français}
%lombardi@math.univ-toulouse.fr
\professor{E.Lombardie}
\usepackage{graphicx}%----pour mettre des images
\usepackage[utf8]{inputenc}%---encodage
\usepackage{geometry}%---pour modifier les tailles et mettre a4paper
%\usepackage{awesomebox}%---pour les boites d'exercices, de pbq et de croquis ---d\'esactiv\'e pour les TP de PC
\usepackage{tikz}%---pour deiffner + d\'ependance de chemfig
% \usepackage{tabularx}%---pour dimensionner automatiquement les tableaux avec variable X
\usepackage{awesomebox}%---Pour les boites info, danger et autres
\usepackage{menukeys}%---Pour deiffner les touches de Calculatrice
\usepackage{fancyhdr}%---pour les en-t\^ete personnalis\'ees
\usepackage{blindtext}%---pour les liens
\usepackage{hyperref}%---pour les liens (\`a mettre en dernier)
\usepackage{caption}%---pour la francisation de la l\'egende table vers Tableau
\usepackage{pifont}
\usepackage{array}%---pour les tableaux
\usepackage{lipsum}
\usepackage{yFlatTable}
\usepackage{multicol}
\newcommand{\Lim}[1]{\lim\limits_{\substack{#1}}\:}
\renewcommand{\vec}{\overrightarrow}
\newcommand{\N}[0]{\mathbb{N}}
\newcommand{\R}[0]{\mathbb{R}}
\newcommand{\dd}{\mathrm{d}}
\newcommand{\norm}[1]{||\vec{#1}||}
\newcommand{\fo}{\psi(\vec{r},t)}
\newcommand{\foe}{\psi(\vec{r},t)\*}
\newcommand{\HH}{\hat{H}}
\newcommand{\hb}{\hbar}
\newcommand{\lap}{\nabla^2}
\newcommand{\lapcc}{\frac{\partial^2 }{\partial x^2}+\frac{\partial^2 }{\partial y^2}+\frac{\partial^2 }{\partial z^2}}
\newcommand{\E}{\vec{E}(M)}
\newcommand{\B}{\vec{B}(M)}
\newcommand{\V}{V(M)}
\newcommand{\er}{\vec{e_{\rho}}}
\newcommand{\ep}{\vec{e_{\varphi}}}
\newcommand{\pip}{\pi^+}
\newcommand{\pim}{\pi^-}
\newcommand{\mv}{\mathcal{V}}
\DeclareMathOperator{\Div}{Div}
\newcommand{\rot}{\vec{\mathrm{rot}}}
\newcommand{\grad}{\vec{\mathrm{grad}}}
\begin{document}
\setcounter{chapter}{1}
%C2 : Applications linéaires sur espaces euclidiens
%C3 : Espaces hermitiens
%C4 : Formes quadratiques
%C5 : Introduction aux groupes

%TD = U-111/U-112
%CM : B-08/B-07/Grignard

%CC2 : TD du 19 mars
%CC1 : 1h30
%CC2 : TD du 7 mai
%CC3 : 2 juin
%CC4 : 9 juin
\chapter{Endomorphismes d'espaces euclidiens}

\section{Rappels sur les projections et symétries}
Soient $F$ et $G$ deux sous-espaces supplémentaires de $E$ ($E = F \oplus G$). Tout vecteur $x \in E$ se décompose de manière unique en $x = x_F + x_G$ avec $x_F \in F$ et $x_G \in G$.

\begin{itemize}
    \item Projection sur $F$ parallèlement à $G$ : $p_{F//G}(x) = x_F$.
    \item Symétrie par rapport à $F$ parallèlement à $G$ : $s_{F//G}(x) = x_F - x_G$.
\end{itemize}

\checkInfo{Lien entre projecteur et symétrie}{On a la relation $s = 2p - Id$, soit $s(x) = 2x_F - (x_F + x_G) = x_F - x_G$.}

\section{Isométries et matrices orthogonales}

\subsection{Définitions et propriétés}

\begin{definition}[Isométrie vectorielle]
    Un endomorphisme $f \in \mathcal{L}(E)$ est une isométrie (ou un endomorphisme orthogonal) si :
    $$\forall x \in E, \|f(x)\| = \|x\|$$
\end{definition}

\begin{proposition}[Caractérisations d'une isométrie]
    Soit $f \in \mathcal{L}(E)$. Les assertions suivantes sont équivalentes :
    \begin{enumerate}
        \item $f$ est une isométrie.
        \item $\forall (x,y) \in E^2, (f(x)|f(y)) = (x|y)$ (conservation du produit scalaire).
        \item L'image de toute base orthonormée (BON) par $f$ est encore une BON.
    \end{enumerate}
\end{proposition}

\begin{definition}[Matrice orthogonale]
    Une matrice $Q \in \mathcal{M}_n(\mathbb{R})$ est dite orthogonale si elle est inversible et que :
    $$Q^{-1} = {}^tQ \iff {}^tQQ = I_n$$
\end{definition}

\begin{proposition}[Matrice d'une isométrie en BON]
    Soit $E$ un espace euclidien. $f$ est une isométrie si et seulement si sa matrice dans une BON est orthogonale.
\end{proposition}

\criticalInfo{Condition cruciale}{Cette équivalence est fausse si la base utilisée n'est pas orthonormée.}

\begin{proposition}[Déterminant d'une matrice orthogonale]
    Si $Q \in O(n)$, alors $\det(Q) = \pm 1$. 
    De plus, pour toute matrice $A \in \mathcal{M}_n(\mathbb{R})$ et tout scalaire $\lambda$ :
    $$\det(\lambda A) = \lambda^n \det(A)$$
\end{proposition}

\begin{proposition}[Valeurs propres réelles d'une isométrie]
    Si $f$ est une isométrie, ses valeurs propres réelles éventuelles appartiennent à $\{-1, 1\}$.
\end{proposition}

\begin{definition}[Groupes orthogonaux]
    On note $O(n)$ le groupe orthogonal (matrices orthogonales) et $SO(n)$ le groupe spécial orthogonal (matrices de rotation) défini par :
    $$SO(n) = \{ Q \in O(n) \mid \det(Q) = 1 \}$$
\end{definition}

\subsection{Étude de O(2)}

\begin{proposition}[Classification des éléments de $O(2)$]
    Soit $A \in O(2)$.
    \begin{itemize}
        \item Si $A \in SO(2)$, $A$ est une rotation d'angle $\theta$ : 
        $A = \begin{pmatrix} \cos \theta & -\sin \theta \\ \sin \theta & \cos \theta \end{pmatrix}$
        \item Si $A \notin SO(2)$ (donc $\det A = -1$), $A$ est une réflexion (symétrie orthogonale par rapport à une droite) :
        $A = \begin{pmatrix} \cos \theta & \sin \theta \\ \sin \theta & -\cos \theta \end{pmatrix}$
    \end{itemize}
\end{proposition}

\subsection{Étude de O(3)}

\begin{proposition}[Forme réduite des éléments de $O(3)$]
    Pour tout $A \in O(3)$, il existe une BON dans laquelle la matrice de $A$ s'écrit :
    $$A' = \begin{pmatrix} \cos \theta & -\sin \theta & 0 \\ \sin \theta & \cos \theta & 0 \\ 0 & 0 & \varepsilon \end{pmatrix}$$
    avec $\varepsilon = \det(A) = \pm 1$.
\end{proposition}

\begin{proposition}[Classification des isométries de $O(3)$]
    Soit $A \in O(3)$ différente de $I_d$.
    \begin{itemize}
        \item Si $\det A = 1$ : $A$ est une rotation d'axe $E_1(A) = \ker(A - I)$. L'angle $\theta$ vérifie $Tr(A) = 1 + 2\cos \theta$.
        \item Si $\det A = -1$ : $A$ est une roto-réflexion. L'angle $\theta$ de la rotation associée vérifie $Tr(A) = -1 + 2\cos \theta$.
    \end{itemize}
\end{proposition}

\section{Adjoint d'un endomorphisme}

\begin{definition}[Adjoint]
    Soit $E$ un espace euclidien et $f \in \mathcal{L}(E)$. Il existe un unique endomorphisme $f^*$, appelé adjoint de $f$, tel que :
    $$\forall (x,y) \in E^2, (f(x)|y) = (x|f^*(y))$$
    En base orthonormée, si $M = Mat(f)$, alors $Mat(f^*) = {}^tM$.
\end{definition}

\begin{proposition}[Propriétés de l'adjoint]
    \begin{itemize}
        \item $(f^*)^* = f$
        \item $(f \circ g)^* = g^* \circ f^*$
        \item Si $f$ est inversible, $(f^*)^{-1} = (f^{-1})^*$
        \item $\ker f^* = (\text{Im } f)^\perp$ et $\text{Im } f^* = (\ker f)^\perp$
    \end{itemize}
\end{proposition}

\section{Endomorphismes auto-adjoints (Symétriques)}

\begin{definition}[Adjoint et matrice]
    $f$ est auto-adjoint (ou symétrique) si $f = f^*$. En BON, cela revient à dire que sa matrice $A$ est symétrique ($A = {}^tA$).
\end{definition}

\begin{theorem}[Théorème Spectral]
    Si $f$ est un endomorphisme auto-adjoint d'un espace euclidien $E$, alors :
    \begin{enumerate}
        \item Toutes les valeurs propres de $f$ sont réelles.
        \item $f$ est diagonalisable dans une base orthonormée de vecteurs propres.
        \item Les sous-espaces propres sont deux à deux orthogonaux.
    \end{enumerate}
\end{theorem}

\checkInfo{Méthode de diagonalisation en BON}{
    1. Vérifier que $A$ est symétrique.
    2. Calculer le polynôme caractéristique et les valeurs propres.
    3. Pour chaque sous-espace propre $E_\lambda$ :
       \begin{itemize}
           \item Déterminer une base.
           \item Si $\dim E_\lambda > 1$, utiliser le procédé de Gram-Schmidt pour la rendre orthonormée.
       \end{itemize}
    4. Concaténer les bases pour obtenir la BON globale.
}

\section{Décomposition polaire et valeurs singulières}

\subsection{Positivité}
\begin{definition}[Positivité]
    $f$ auto-adjoint est positif si $\forall x \in E, (x|f(x)) \ge 0$.
    Il est défini positif si $\forall x \neq 0, (x|f(x)) > 0$.
\end{definition}

\begin{proposition}[Critère spectral de positivité]
    $f$ est positif ssi ses valeurs propres sont dans $\mathbb{R}^+$.
    $f$ est défini positif ssi ses valeurs propres sont dans $\mathbb{R}^{+*}$.
\end{proposition}

\subsection{Décomposition polaire}
\begin{theorem}
    Pour toute matrice $M \in GL_n(\mathbb{R})$, il existe un unique couple $(Q, S)$ avec $Q \in O(n)$ et $S$ symétrique définie positive tel que $M = QS$.
\end{theorem}

\subsection{Valeurs singulières (SVD)}
\begin{lemma}
    Soit $A \in \mathcal{M}_{n,m}(\mathbb{R})$. La matrice ${}^tAA$ est une matrice symétrique positive de taille $m \times m$.
\end{lemma}

\begin{theorem}[SVD]
    Soit $A \in \mathcal{M}_{n,m}(\mathbb{R})$ de rang $r$. Il existe $Q \in O(n)$, $P \in O(m)$ et des réels $\sigma_1 \ge \sigma_2 \ge \dots \ge \sigma_r > 0$ (valeurs singulières) tels que :
    $$A = Q \Sigma {}^tP$$
    où $\Sigma$ est une matrice de même taille que $A$ contenant les $\sigma_i$ sur sa "diagonale".
\end{theorem}

\tipsInfo{Lien valeurs propres/singulières}{Les valeurs singulières $\sigma_i$ de $A$ sont les racines carrées des valeurs propres non nulles de ${}^tAA$.}

\checkInfo{Méthode pour décomposer en valeurs singulières}{
\begin{enumerate}
    \item Calculer $^tAA$ : on s'assure que c'est symétrique
    \item Diagonaliser $^tAA$ en bon en rangeant les valeurs propres de $^tAA$ par ordre décroissant.
    \item On pose $\sigma_k = \sqrt{\lambda_k}, q_k = \frac{1}{\sigma_k}A_{pk}$
    \item Le théorème assure que la famille orthonormale complète cette famille. On complète cette famille pour obtenir une bon de R en utiisant Gram Schmitch dans $Vect(q_1,\dots, q_r)^\perp$.
\end{enumerate}

}

\end{document}